\chapter{ฟิสิกส์จลนศาสตร์และการประมวลผลสัญญาณตรวจจับการก้าว}
\label{chap:ch02}

\section*{แผนบริหารการสอนประจำบทที่ 2}
\addcontentsline{toc}{section}{แผนบริหารการสอนประจำบทที่ 2}

\noindent\textbf{หัวข้อเนื้อหาประจำบท}
\begin{enumerate}
    \item จลนศาสตร์การเดินของมนุษย์และวัฏจักรการก้าว
    \item การคำนวณเวกเตอร์ขนาดความเร่งสามมิติแบบอิสระต่อการจัดวางพิกัด
    \item การแยกเวกเตอร์ความโน้มถ่วงของโลกและการกรองความถี่ต่ำ
    \item อัลกอริทึมตรวจจับยอดพีคแบบปรับเกณฑ์พลวัต
    \item กลไกระยะพักฟื้นทางชีวกลศาสตร์เพื่อป้องกันการนับก้าวซ้ำ
    \item การคัดกรองสัญญาณรบกวนจากการสั่นสะเทือนของยานพาหนะและการแกว่งแขนหลอก
\end{enumerate}

\noindent\textbf{วัตถุประสงค์เชิงพฤติกรรม}
\begin{enumerate}
    \item อธิบายวัฏจักรการก้าวเดินของมนุษย์ในแต่ละระยะจลนศาสตร์ได้อย่างถูกต้อง
    \item พิสูจน์สมการเวกเตอร์ขนาดรวมความเร่งแบบยูคลิดสามมิติได้อย่างถูกต้อง
    \item ออกแบบและวิเคราะห์การทำงานของอัลกอริทึม Dynamic Peak Detection พร้อมเกณฑ์ปรับตัวได้อย่างถูกต้อง
    \item ประยุกต์ใช้ค่าจำกัดระยะพักฟื้นทางสรีรวิทยา $250\text{ ms}$ ในการกำจัดสัญญาณรบกวนได้อย่างมีประสิทธิภาพ
\end{enumerate}

\noindent\textbf{กิจกรรมการเรียนการสอน}
\begin{enumerate}
    \item บรรยายทฤษฎีจลนศาสตร์การเคลื่อนที่และแรงกระทำต่อข้อต่อในขณะก้าวเดิน
    \item การสาธิตการคำนวณเวกเตอร์ลัพธ์ความเร่งสามมิติด้วยฟังก์ชันคณิตศาสตร์บนกระดาน
    \item ปฏิบัติการจำลองสัญญาณคลื่นความเร่งและการปรับระดับเส้นเกณฑ์ตัดยอดพีคบนโปรแกรมคอมพิวเตอร์
    \item การสรุปบทเรียนและทำแบบฝึกหัดคำถามทบทวนท้ายบท
\end{enumerate}

\noindent\textbf{สื่อการเรียนการสอน}
\begin{enumerate}
    \item เอกสารคำสอนบทที่ 2 เรื่องฟิสิกส์จลนศาสตร์และการประมวลผลสัญญาณ
    \item โปรแกรมจำลองกราฟคลื่นสัญญาณความเร่งดิจิทัล
    \item ชุดข้อมูลสัญญาณความเร่งจากการเดิน การวิ่ง และการนั่งยานพาหนะจริง
\end{enumerate}

\noindent\textbf{การประเมินผล}
\begin{enumerate}
    \item การประเมินความเข้าใจผ่านการตอบคำถามในชั้นเรียน
    \item การตรวจความถูกต้องของแบบฝึกหัดคำนวณค่าเวกเตอร์และเกณฑ์พลวัต
    \item การประเมินผลจากคำถามทบทวนท้ายบทที่ 2
\end{enumerate}

\newpage

\section{จลนศาสตร์การเดินของมนุษย์และวัฏจักรการก้าว}
การเดินของมนุษย์ (Human Gait) เป็นกระบวนการเคลื่อนที่ของจุดศูนย์ถ่วงร่างกาย (Center of Mass: COM) ไปข้างหน้าอย่างต่อเนื่อง วัฏจักรการก้าวเดิน (Gait Cycle) ของขาข้างหนึ่งจะเริ่มต้นขึ้นเมื่อส้นเท้าสัมผัสพื้น (Heel Strike) จนกระทั่งส้นเท้าของขาข้างเดิมสัมผัสพื้นอีกครั้ง ซึ่งแบ่งออกเป็น 2 ช่วงจลนศาสตร์สำคัญ

\begin{enumerate}
    \item \textbf{ช่วงเท้าแตะพื้น (Stance Phase)} คิดเป็นประมาณร้อยละ 60 ของวัฏจักรการก้าว ประกอบด้วยจังหวะที่ส้นเท้ากระแทกพื้น (Initial Contact) ฝ่าเท้าราบกับพื้น (Foot Flat) การรับน้ำหนักสูงสุด (Midstance) และการดันปลายเท้าถีบตัวออกจากพื้น (Toe-off) ในช่วงนี้จะเกิดแรงปฏิกิริยาจากพื้นสู่ร่างกาย (Ground Reaction Force: GRF) ส่งผ่านกระดูกสันหลังสู่จุดที่สมาร์ทโฟนแนบอยู่ เกิดเป็นคลื่นความเร่งพุ่งขึ้นเป็นยอดแหลม
    \item \textbf{ช่วงเท้าลอยพ้นพื้น (Swing Phase)} คิดเป็นประมาณร้อยละ 40 ของวัฏจักรการก้าว เป็นจังหวะที่เท้าถูกยกเหวี่ยงไปข้างหน้าเพื่อเตรียมแตะพื้นในก้าวถัดไป ความเร่งสุทธิของร่างกายจะลดลงสู่ระดับต่ำสุด
\end{enumerate}

\section{การคำนวณเวกเตอร์ขนาดความเร่งรวมสามมิติ}
ปัญหาคลาสสิกของแอปพลิเคชันนับก้าวบนสมาร์ทโฟน คือสมาร์ทโฟนสามารถถูกจัดวางได้ในหลากหลายอิริยาบถ เช่น เสียบในกระเป๋ากางเกงแนวนอน กระเป๋าเสื้อแนวตั้ง หรือถือไว้ในมือขณะเดิน ทำให้แกนอ้างอิง $x, y, z$ ของตัวเซนเซอร์หมุนเปลี่ยนทิศทางตลอดเวลา

เพื่อแก้ปัญหานี้อย่างเป็นระบบตามหลักฟิสิกส์ ระบบ Smart Health Tracker จึงใช้วิธีการคำนวณขนาดเวกเตอร์ความเร่งรวมแบบยูคลิดสามมิติ (3D Euclidean Acceleration Magnitude Norm, $a_{\text{norm}}$) ซึ่งมีสมบัติคงรูปไม่แปรผันตามการหมุนของพิกัดอ้างอิง (Rotation-Invariant)

\begin{theorybox}
\textbf{ทฤษฎีเวกเตอร์ขนาดความเร่งรวมอิสระจากการหมุน (Rotation Invariance)}

กำหนดให้พิกัดความเร่งตามแกนของเซนเซอร์คือ $\vec{a} = a_x\hat{i} + a_y\hat{j} + a_z\hat{k}$ ขนาดสเกลาร์รวมของความเร่งคำนวณได้จากผลคูณภายใน (Inner Product) ดังสมการ
\begin{equation}
a_{\text{norm}} = \|\vec{a}\| = \sqrt{a_x^2 + a_y^2 + a_z^2}
\end{equation}
ไม่ว่าสมาร์ทโฟนจะเอียงหรือหมุนรอบแกนใดด้วยเมทริกซ์การหมุน $\mathbf{R} \in SO(3)$ เวกเตอร์ความเร่งใหม่ $\vec{a}' = \mathbf{R}\vec{a}$ จะยังมีขนาดเท่าเดิมเสมอเนื่องจากสมบัติออร์โทโกนัลของเมทริกซ์การหมุน
\begin{equation}
\|\vec{a}'\|^2 = (\mathbf{R}\vec{a})^T (\mathbf{R}\vec{a}) = \vec{a}^T \mathbf{R}^T \mathbf{R} \vec{a} = \vec{a}^T \mathbf{I} \vec{a} = \|\vec{a}\|^2
\end{equation}
ในสภาวะที่ผู้ใช้อยู่นิ่ง ค่า $a_{\text{norm}}$ จะมีค่าเท่ากับขนาดของความเร่งเนื่องจากแรงโน้มถ่วงโลก คือ $g \approx 9.80665\text{ m/s}^2$ พอดี
\end{theorybox}

\section{การสกัดเวกเตอร์ความโน้มถ่วงและการปรับระดับฐาน}
เนื่องจากค่า $a_{\text{norm}}$ ประกอบด้วยค่าคงที่ของแรงโน้มถ่วงโลก $g$ บวกกับความเร่งเชิงพลวัตจากการเคลื่อนไหวร่างกาย ($a_{\text{dynamic}}$) ระบบจึงต้องทำการลบค่าแรงโน้มถ่วงออกเพื่อหาความเร่งสุทธิที่แท้จริง
\begin{equation}
a_{\text{dynamic}}(t) = |a_{\text{norm}}(t) - g_{\text{est}}|
\end{equation}
โดยที่ $g_{\text{est}}$ คือค่าแรงโน้มถ่วงที่ประเมินผ่านตัวกรองความถี่ต่ำผ่าน (Low-Pass Filter) เพื่อรองรับกรณีที่เซนเซอร์ในสมาร์ทโฟนบางรุ่นอาจมีการคลาดเคลื่อนของจุดกำเนิดศูนย์ (Zero-g Offset)

\section{อัลกอริทึม Dynamic Peak Detection}
การใช้เกณฑ์การตัดสินใจแบบค่าคงที่ (Static Threshold) เช่น ตัดสินว่าก้าวเมื่อ $a_{\text{dynamic}} > 2.0\text{ m/s}^2$ มีข้อจำกัดร้ายแรง เนื่องจากแต่ละบุคคลมีแรงกระแทกในการเดินไม่เท่ากัน และความเร็วในการเดินหรือวิ่งจะทำให้แอมพลิจูดของคลื่นเปลี่ยนไปตลอดเวลา

ระบบจึงเลือกใช้อัลกอริทึมการตรวจจับยอดพีคแบบปรับเกณฑ์พลวัต (Dynamic Peak Detection Algorithm) ซึ่งจะคำนวณค่าเฉลี่ยเคลื่อนที่และค่าเบี่ยงเบนมาตรฐานของสัญญาณแบบเรียลไทม์

\begin{equation}
T_{\text{dynamic}}(t) = \mu_{w}(t) + \beta \cdot \sigma_{w}(t)
\end{equation}
โดยที่ $\mu_w(t)$ คือค่าเฉลี่ยของขนาดความเร่งในหน้าต่างเวลา $w$ ขนาด 50 ตัวอย่าง ($1$ วินาที) $\sigma_w(t)$ คือค่าเบี่ยงเบนมาตรฐานของการแกว่งตัว และ $\beta$ คือตัวคูณปรับความไว (Sensitivity Factor โดยทั่วไปกำหนดที่ $0.65 - 0.85$) เมื่อสัญญาณ $a_{\text{dynamic}}(t)$ พุ่งทะลุ $T_{\text{dynamic}}(t)$ และมีอนุพันธ์เปลี่ยนเครื่องหมายจากบวกเป็นลบ ($\frac{da}{dt} = 0$) จุดดังกล่าวจะถูกระบุเป็นยอดพีคของก้าวเดิน

\begin{figure}[H]
\centering
\begin{tikzpicture}[scale=0.92]
  % Grid
  \draw[step=1cm, gray!20, very thin] (0, 0) grid (12, 5);

  % Axes
  \draw[->, >=stealth, line width=1pt] (0, 0) -- (12.5, 0) node[right] {\small เวลา (วินาที)};
  \draw[->, >=stealth, line width=1pt] (0, 0) -- (0, 5.2) node[above] {\small $a_{\text{norm}}\ (\text{m/s}^2)$};

  % Gravity Baseline
  \draw[dashed, color=gray!80, line width=0.8pt] (0, 1.8) -- (12, 1.8) node[right] {\scriptsize $g \approx 9.8\text{ m/s}^2$};

  % Dynamic Threshold Line
  \draw[color=rbruGold, line width=1.2pt] plot[smooth] coordinates {
    (0, 2.6) (2, 2.7) (4, 2.9) (6, 2.8) (8, 3.1) (10, 2.8) (12, 2.7)
  };
  \node[color=rbruGold, anchor=west] at (10.2, 3.4) {\scriptsize เกณฑ์พลวัต $T_{\text{dynamic}}$};

  % Acceleration Signal Wave
  \draw[color=rbruNavy, line width=1.4pt] plot[smooth] coordinates {
    (0, 1.8) (0.8, 1.4) (1.5, 3.8) (2.2, 1.2) (2.8, 1.8)
    (3.5, 1.3) (4.2, 4.4) (4.9, 1.1) (5.5, 1.7)
    (6.3, 1.5) (7.0, 4.2) (7.7, 1.2) (8.4, 1.9)
    (9.1, 1.4) (9.8, 4.6) (10.5, 1.0) (11.2, 1.8) (12, 1.8)
  };

  % Step Peak Indicators
  \fill[color=rbruEmerald] (1.5, 3.8) circle (3.5pt) node[above=2pt] {\scriptsize ก้าวที่ 1};
  \fill[color=rbruEmerald] (4.2, 4.4) circle (3.5pt) node[above=2pt] {\scriptsize ก้าวที่ 2};
  \fill[color=rbruEmerald] (7.0, 4.2) circle (3.5pt) node[above=2pt] {\scriptsize ก้าวที่ 3};
  \fill[color=rbruEmerald] (9.8, 4.6) circle (3.5pt) node[above=2pt] {\scriptsize ก้าวที่ 4};

  % Refractory Window visualization
  \draw[fill=rbruCyan!20, opacity=0.7] (1.5, 0.2) rectangle (3.0, 0.6);
  \node at (2.25, 0.4) {\tiny $T_{\text{refract}} = 250\text{ ms}$};
  \draw[<->, >=stealth, line width=0.6pt] (1.5, 0.4) -- (3.0, 0.4);
\end{tikzpicture}
\caption{รูปคลื่นความเร่งสามมิติ ยอดพีคของก้าวเดิน และช่วงระยะพักฟื้นชีวภาพ $T_{\text{refract}}$}
\label{fig:step_waveform_refractory}
\end{figure}

\section{กลไกระยะพักฟื้นทางชีวกลศาสตร์ (Refractory Period)}
ในทางชีวกลศาสตร์ของมนุษย์ ขีดจำกัดทางสรีรวิทยาของการเคลื่อนไหวก้าวขาที่เร็วที่สุด (เช่น นักวิ่งระยะสั้นระดับโลก 100 เมตร) จะมีอัตราความถี่ก้าวสูงสุดไม่เกิน 4 ถึง 5 ก้าวต่อวินาที นั่นหมายความว่า คาบเวลาระหว่างก้าวสองก้าวที่ติดกัน ($t_{k+1} - t_k$) จะไม่มีทางน้อยกว่า $200$ ถึง $250\text{ มิลลิวินาที}$ ได้อย่างแน่นอน

ระบบจึงวางเงื่อนไขระยะพักฟื้นทางชีวกลศาสตร์ (Refractory Period Guard, $T_{\text{refract}} = 250\text{ ms}$) เมื่อตรวจพบยอดพีคที่ผ่านเกณฑ์แล้ว ตัวนับเวลาจะล็อกไม่ให้ยอมรับยอดพีคอื่นใดที่เกิดขึ้นภายในระยะเวลา $250\text{ ms}$ ต่อจากนั้น แม้ว่าสัญญาณจะมีความสูงทะลุเกณฑ์ก็ตาม กลไกนี้มีประสิทธิภาพสูงยิ่งในการตัดสัญญาณคลื่นสะท้อนรอง (Secondary Bounces) ที่เกิดจากการสั่นของกล้ามเนื้อหรือการสะเทือนของเนื้อผ้าในกระเป๋า

\section{การคัดกรองสัญญาณรบกวนจากการสั่นสะเทือน}
ข้อแตกต่างที่สำคัญระหว่างการก้าวเดินจริงกับการนั่งบนยานพาหนะ (รถยนต์ มอเตอร์ไซค์ หรือรถไฟ) คือลักษณะของรูปแบบสเปกตรัมความถี่
\begin{enumerate}
    \item \textbf{การสั่นสะเทือนของยานพาหนะ} จะมีความถี่สูงกว่า $10 - 50\text{ Hz}$ มีแอมพลิจูดสม่ำเสมอแต่ไม่มีจังหวะการเปลี่ยนเฟสแบบก้าวขา และไม่มีระยะสวิงของร่างกาย
    \item \textbf{การแกว่งแขนขณะนั่งอยู่กับที่} มีแอมพลิจูดเกิดขึ้นเพียงช่วงสั้นๆ ไม่เกิน $2 - 3$ ครั้งแล้วหยุดนิ่ง
\end{enumerate}

ระบบจึงได้เสริมเงื่อนไขการตรวจสอบความต่อเนื่องของก้าว (Step Continuity Verification) โดยจะเริ่มบันทึกก้าวเข้าสู่หน่วยความจำสะสมหลังจากตรวจพบก้าวที่สอดคล้องตามเกณฑ์ติดต่อกันอย่างน้อย 4 ถึง 5 ก้าว หากหยุดชะงักก่อนหน้านั้น ก้าวหลอกดังกล่าวจะถูกตัดทิ้งเป็นสัญญาณรบกวน

\section{สถาปัตยกรรมโครงข่ายประสาทเทียมเชิงลึกและการเรียนรู้ระดับอุปกรณ์ (On-Device Deep Learning)}
เพื่อยกระดับการตรวจวัดสุขภาพจากการนับก้าวเชิงตัวเลขแบบธรรมดา สู่การวินิจฉัยเชิงชีวกลศาสตร์ขั้นสูง ระบบ Smart Health Tracker ได้บรรจุโมเดลโครงข่ายประสาทเทียมเชิงลึกหลายชั้น (Deep Multilayer Perceptron: MLP) ที่ทำงานและฝึกสอนได้จริงบนอุปกรณ์ของผู้ใช้ (On-Device Continual Deep Learning)

\begin{theorybox}
\textbf{แบบจำลองทางคณิตศาสตร์และการส่งสัญญาณไปข้างหน้า (Forward Propagation)}

โมเดลรับเวกเตอร์คุณลักษณะจลนศาสตร์ 6 มิติ $\mathbf{x} = [x_1, x_2, x_3, x_4, x_5, x_6]^T$ ซึ่งสกัดจากข้อมูลรอบ 4 ชั่วโมง (จำนวนก้าว, จังหวะความถี่ก้าวเฉลี่ย, ความเข้มข้นของการเคลื่อนไหว, สัดส่วนเวลาการนั่งนิ่ง, การเผาผลาญพลังงาน, และอัตราการสูญเสียน้ำ)

โครงข่ายประกอบด้วย 3 ชั้น ได้แก่ ชั้นซ่อนที่หนึ่ง (8 โหนด, ฟังก์ชันกระตุ้น LeakyReLU), ชั้นซ่อนที่สอง (6 โหนด, ฟังก์ชันกระตุ้น Tanh), และชั้นส่งออก (4 โหนด, ฟังก์ชันกระตุ้น Softmax)
\begin{align}
\mathbf{h}_1 &= \text{LeakyReLU}(\mathbf{W}_1 \mathbf{x} + \mathbf{b}_1) \\
\mathbf{h}_2 &= \tanh(\mathbf{W}_2 \mathbf{h}_1 + \mathbf{b}_2) \\
\hat{\mathbf{y}} &= \text{Softmax}(\mathbf{W}_3 \mathbf{h}_2 + \mathbf{b}_3)
\end{align}
โดยที่ $\hat{\mathbf{y}} = [\hat{y}_1, \hat{y}_2, \hat{y}_3, \hat{y}_4]^T$ แทนความน่าจะเป็นของ 4 สภาวะทางชีวกลศาสตร์ ได้แก่ ก้าวปกติสมดุล (Balanced Gait), กล้ามเนื้อล้าสะสม (Muscle Fatigue), แอโรบิกประสิทธิภาพสูง (Aerobic), และภาวะหลอดเลือดดำชะงักจากการนั่งนาน (Sedentary Risk)
\end{theorybox}

เมื่อผู้ใช้บันทึกข้อมูลการเคลื่อนไหวใหม่ ระบบสามารถกระตุ้นกระบวนการฝึกสอนย้อนกลับ (Backpropagation) เพื่อปรับค่าน้ำหนัก $\mathbf{W}$ และไบแอส $\mathbf{b}$ ให้สอดคล้องกับชีวกลศาสตร์เฉพาะบุคคลตามฟังก์ชันการสูญเสีย Cross-Entropy Loss
\begin{equation}
\mathcal{L} = -\sum_{c=1}^4 y_c \ln(\hat{y}_c)
\end{equation}
และทำการปรับปรุงค่าพารามิเตอร์ผ่านเกรเดียนต์ดิเซนต์แบบสุ่ม (Stochastic Gradient Descent: SGD) ด้วยอัตราการเรียนรู้ $\eta$
\begin{equation}
\mathbf{W}_l^{(t+1)} = \mathbf{W}_l^{(t)} - \eta \frac{\partial \mathcal{L}}{\partial \mathbf{W}_l}, \quad \mathbf{b}_l^{(t+1)} = \mathbf{b}_l^{(t)} - \eta \frac{\partial \mathcal{L}}{\partial \mathbf{b}_l}
\end{equation}
กระบวนการนี้ทำให้แอปพลิเคชันสามารถเรียนรู้และปรับตัวเข้ากับรูปแบบการเดินเฉพาะบุคคลได้อย่างแท้จริง โดยไม่ต้องพึ่งพาเซิร์ฟเวอร์ภายนอก

\section{คณิตศาสตร์จลนศาสตร์พิกัดภูมิศาสตร์และสมการฮาเวอร์ซีน}
ในการคำนวณระยะทางจริงจากการเคลื่อนที่บนพื้นผิวพิภพที่มีลักษณะเป็นทรงกลมแป้น (Oblate Spheroid) การใช้สูตรระยะทางยูคลิดบนระนาบ 2 มิติ ($d = \sqrt{\Delta x^2 + \Delta y^2}$) จะทำให้เกิดความคลาดเคลื่อนทางเรขาคณิตทรงกลมอย่างรุนแรง

ระบบ Smart Health Tracker จึงประยุกต์ใช้ \textbf{สมการฮาเวอร์ซีน (Haversine Formula)} เพื่อคำนวณระยะทางบนผิวโค้งวงกลมใหญ่ (Great-Circle Distance) ระหว่างสองจุดพิกัด $(\phi_1, \lambda_1)$ และ $(\phi_2, \lambda_2)$

\begin{theorybox}
\textbf{สมการการคำนวณระยะทางผิวโค้งโลกของฮาเวอร์ซีน}

กำหนดให้ $\phi$ คือละติจูดในหน่วยเรเดียน, $\lambda$ คือลองจิจูดในหน่วยเรเดียน และ $R = 6{,}371{,}000\text{ เมตร}$ คือรัศมีเฉลี่ยของโลก
\begin{align}
\Delta \phi &= \phi_2 - \phi_1 \\
\Delta \lambda &= \lambda_2 - \lambda_1 \\
a &= \sin^2\left(\frac{\Delta \phi}{2}\right) + \cos(\phi_1)\cos(\phi_2)\sin^2\left(\frac{\Delta \lambda}{2}\right) \\
c &= 2 \operatorname{atan2}\left(\sqrt{a}, \sqrt{1-a}\right) \\
d &= R \times c
\end{align}
ความเร็วขณะใดขณะหนึ่ง ($v_{\text{inst}}$) คำนวณได้จากผลต่างระยะทางต่อผลต่างเวลา
\begin{equation}
v_{\text{inst}} = \frac{d}{\Delta t} \times 3.6\text{ กิโลเมตรต่อชั่วโมง}
\end{equation}
และทิศทางการเคลื่อนที่ (Bearing หรือ Azimuth, $\theta$) บนระนาบราบคำนวณจาก
\begin{equation}
\theta = \operatorname{atan2}\left(\sin(\Delta \lambda)\cos(\phi_2), \; \cos(\phi_1)\sin(\phi_2) - \sin(\phi_1)\cos(\phi_2)\cos(\Delta \lambda)\right)
\end{equation}
\end{theorybox}

เพื่อขจัดปัญหาการแกว่งตัวของพิกัดขณะผู้ใช้หยุดยืนอยู่กับที่ (Stationary GPS Jitter) และสัญญาณสะท้อนจากอาคารสูง (Multipath Interference) ระบบได้กำหนดฟังก์ชันกรองสัญญาณเชิงตรรกะ 3 ด่าน
\begin{equation}
\text{Accept}(P_k) = 
\begin{cases}
\text{True} & \text{ถ้า } \text{Accuracy}(P_k) \le 25\text{ m} \ \land \ d(P_{k-1}, P_k) \ge 2.0\text{ m} \ \land \ v_{\text{inst}} \le 45\text{ km/h} \\
\text{False} & \text{กรณีอื่นๆ (ตัดทิ้งเป็นสัญญาณรบกวน)}
\end{cases}
\end{equation}

\section*{คำถามทบทวนท้ายบทที่ 2}
\addcontentsline{toc}{section}{คำถามทบทวนท้ายบทที่ 2}
\begin{enumerate}
    \item จงอธิบายความแตกต่างระหว่างระยะ Stance Phase และ Swing Phase ในแง่ของแรงปฏิกิริยาจากพื้น
    \item เหตุใดเวกเตอร์ขนาดความเร่งแบบยูคลิด $a_{\text{norm}}$ จึงไม่แปรผันตามการหมุนของโทรศัพท์ จงแสดงการพิสูจน์
    \item หากใช้เกณฑ์ตัดสินการก้าวเดินแบบคงที่ (Static Threshold) จะเกิดผลกระทบอย่างไรเมื่อผู้ใช้เปลี่ยนจากการเดินช้าเป็นการวิ่งเร็ว
    \item จงอธิบายหลักการทำงานของ Refractory Period $250\text{ ms}$ และประโยชน์ในการลดความผิดพลาดของการนับก้าว
    \item จงเขียนสมการฮาเวอร์ซีนและอธิบายเหตุผลทางฟิสิกส์ที่ต้องใช้สูตรนี้แทนระยะทางแบบยูคลิดในการหาเส้นทางการเดินทางจริง
    \item จงอธิบายกลไกการกรองสัญญาณรบกวนของ GPS เพื่อป้องกันปรากฏการณ์พิกัดดริฟต์ขณะหยุดนิ่ง
\end{enumerate}
\clearpage
